% Nouveau EGA V (cote 25): archive cote 25, site page 10 (PDF page 11).
% typescript. Draft compared: none.
% AI-transcribed by Claude Opus 5.5 (claude-opus-5-5, adaptive thinking, effort high), Message Batches API, batch ega25b (msgbatch_01A7n5YuXubUPNK1zsepi8RG), 2026-11-13; not yet checked by eye.
\scanpage{25/10}
\mspage{7}

On est ramené au cas où X est irréductible, grâce à 2.5. Par construction
même, Y est défini à partir de $X_{k(\hand{\eta})}$ comme le diviseur d'une section
d'un module invers\unsure{i|a}ble sur $X_{k(\hand{\eta})}$ (savoir de l'image inverse de $\underline{O}_P(1)$))
D'autre part, $X_{k(\hand{\eta})}$ est irréductible (car X l'est, et $k(\hand{\eta})$ est une
\margin{|} extension transcendante pure de k, ce qu'on aurait dû signaler au début du
N° ...), et $Y_{\hand{\eta}} \neq X_{k(\hand{\eta})}$ , puisque l'image de $Y_{\hand{\eta}}$ dans X (contrairement
à celle de $X_{k(\hand{\eta})}$ , qui est fid. plat sur X) n'est pas égale à X, ne conte-%
nant pas en effet les points fermés de X en vertu de 2.3. Il s'ensuit que
dim $Y_{\hand{\eta}}$ = dim $X_{k(\hand{\eta})}$ -1 \inserted{\unsure{(n°)|(nf)}} = dim X - 1 , cqfd.

\uline{Proposition }2.7. Soit F un Module quasi-cohérent sur X, d'où $G_{\hand{\eta}}$ sur $Y_{\hand{\eta}}$ .
Soient $Z_i$ les cycles premiers associés à F tels que dim $\hand{\overline{f(Z_i)}}$ $>$ 0. Soit
$Z_{i\hand{\eta}}$ l'image inverse de $Z_i$ dans $Y_{\hand{\eta}}$ , alors les $Z_{i\hand{\eta}}$ sont exactement tous
les cycles premiers associés à $G_{\hand{\eta}}$ , de plus leurs relations d'inclusion
sont celles entre les $Z_i$ .

La dernière assertion est contenue dans 2.4. D'autre part, comme
$\underline{Y} \rightarrow X$ est un fibré projectif, donc plat\xout{b} à fibres $(S_1)$ et irr. , il résul
te du par.3 que les cycles premiers associés à l'image inverse \inserted{G} de F sur
$\overline{Y}$ sont les images inverses \xout{xxxxxxxx} des cycles premiers associés à F,
d'où en \xout{prenant} \inserted[typed]{s'induisant sur} la fibre générique \inserted[typed]{$Y_{\hand{\eta}}$} de $\underline{Y}$ sur $\check{P}$ , le fait que les
cycles premiers associés à $G_{\hand{\eta}}$ sont les images inverses \uline{non vides} des $Z_i$ ,
ce qui établit 2.7. moyennant 2.3. A vrai dire, le passage par $\underline{Y}$ est
inutile, on peut utiliser directement le fait que $Y_{\hand{\eta}}$ \hand{$\longrightarrow$} X est\xout{n} plat à
fibres $(S_1)$ (et même géom. régulières, i.e. le morphisme est régulier)
et à fibres irréductibles (et même géométriquement irréductibles : c'est
des localisés de schémas projectifs)\hand{;} \hand{\unsure{remarque} \unsure{analogue} \unsure{pour} \unsure{la}
dim. de 2.\unsure{1}.}

\margin{10}
